% GENERATED FILE. Edit canonical lessons, then run npm run generate:books.
% canonical-source-hash: fnv1a64:93ee76d1659a8c97
% canonical-lessons: TE-C275-pujari, TE-C275-kalakarudu, TE-C275-kavi, TE-C275-sangitam, TE-C275-natakam

\chapter{Priest, Artist, Poet, Music, Play}
\label{ch:te-priest-artist-poet-music-play}
\hlchaptermodality{275}

\begin{chapteropening}
\textbf{By the end of this chapter:} \emph{I can say a priest, an artist, a poet, music and a play.}

Complete the last lesson of chapter 275: I can say a priest, an artist, a poet, music and a play.
\end{chapteropening}

\section[\texorpdfstring{\te{పూజారి} (pūjāri)}{పూజారి (pūjāri)}]{\te{పూజారి} (pūjāri) — a priest}

\label{lesson:TE-C275-pujari}



\begin{quote}
Before the new one: say the Telugu for a subject, then the Telugu for mathematics.
\end{quote}

\subsection*{You'll want to know: \te{పూజారి}}
\textbf{\te{పూజారి}} — \emph{pūjāri} — \textquotedblleft{}a priest\textquotedblright{}.

\begin{grammarlens}[title={What you've built}]
One more word for this chapter.
\end{grammarlens}

\subsection*{Your turn}
\begin{itemize}
\raggedright
  \item \emph{Say it:} \emph{pūjāri}
  \item \emph{Say it:} \emph{pūjāri}, once more
  \item \emph{Say it:} say \emph{pūjāri}
  \item \emph{Recall:} say the Telugu for a subject, then the Telugu for mathematics, then say \emph{pūjāri} again
\end{itemize}

\begin{tcolorbox}[breakable,colback=teal!4,colframe=teal!35!black,title={Before you move on}]
What does \te{పూజారి} mean? (\textquotedblleft{}A priest\textquotedblright{}.) Say it once more.
\end{tcolorbox}
\section[\texorpdfstring{\te{కళాకారుడు} (kaḷākāruḍu)}{కళాకారుడు (kaḷākāruḍu)}]{\te{కళాకారుడు} (kaḷākāruḍu) — an artist, a performer}

\label{lesson:TE-C275-kalakarudu}



\begin{quote}
Before the new one: say the Telugu for mathematics, then the Telugu for a priest.
\end{quote}

\subsection*{You'll want to know: \te{కళాకారుడు}}
\textbf{\te{కళాకారుడు}} — \emph{kaḷākāruḍu} — \textquotedblleft{}an artist, a performer\textquotedblright{}.

\begin{grammarlens}[title={What you've built}]
One more word for this chapter.
\end{grammarlens}

\subsection*{Your turn}
\begin{itemize}
\raggedright
  \item \emph{Say it:} \emph{kaḷākāruḍu}
  \item \emph{Say it:} \emph{kaḷākāruḍu}, once more
  \item \emph{Say it:} say \emph{kaḷākāruḍu}
  \item \emph{Recall:} say the Telugu for mathematics, then the Telugu for a priest, then say \emph{kaḷākāruḍu} again
\end{itemize}

\begin{tcolorbox}[breakable,colback=teal!4,colframe=teal!35!black,title={Before you move on}]
What does \te{కళాకారుడు} mean? (\textquotedblleft{}An artist, a performer\textquotedblright{}.) Say it once more.
\end{tcolorbox}
\section[\texorpdfstring{\te{కవి} (kavi)}{కవి (kavi)}]{\te{కవి} (kavi) — a poet}

\label{lesson:TE-C275-kavi}



\begin{quote}
Before the new one: say the Telugu for a priest, then the Telugu for an artist.
\end{quote}

\subsection*{You'll want to know: \te{కవి}}
\textbf{\te{కవి}} — \emph{kavi} — \textquotedblleft{}a poet\textquotedblright{}.

\begin{grammarlens}[title={What you've built}]
One more word for this chapter.
\end{grammarlens}

\subsection*{Your turn}
\begin{itemize}
\raggedright
  \item \emph{Say it:} \emph{kavi}
  \item \emph{Say it:} \emph{kavi}, once more
  \item \emph{Say it:} say \emph{kavi}
  \item \emph{Recall:} say the Telugu for a priest, then the Telugu for an artist, then say \emph{kavi} again
\end{itemize}

\begin{tcolorbox}[breakable,colback=teal!4,colframe=teal!35!black,title={Before you move on}]
What does \te{కవి} mean? (\textquotedblleft{}A poet\textquotedblright{}.) Say it once more.
\end{tcolorbox}
\section[\texorpdfstring{\te{సంగీతం} (saṅgītaṁ)}{సంగీతం (saṅgītaṁ)}]{\te{సంగీతం} (saṅgītaṁ) — music}

\label{lesson:TE-C275-sangitam}



\begin{quote}
Before the new one: say the Telugu for an artist, then the Telugu for a poet.
\end{quote}

\subsection*{You'll want to know: \te{సంగీతం}}
\textbf{\te{సంగీతం}} — \emph{saṅgītaṁ} — \textquotedblleft{}music\textquotedblright{}.

\begin{grammarlens}[title={What you've built}]
One more word for this chapter.
\end{grammarlens}

\subsection*{Your turn}
\begin{itemize}
\raggedright
  \item \emph{Say it:} \emph{saṅgītaṁ}
  \item \emph{Say it:} \emph{saṅgītaṁ}, once more
  \item \emph{Say it:} say \emph{saṅgītaṁ}
  \item \emph{Recall:} say the Telugu for an artist, then the Telugu for a poet, then say \emph{saṅgītaṁ} again
\end{itemize}

\begin{tcolorbox}[breakable,colback=teal!4,colframe=teal!35!black,title={Before you move on}]
What does \te{సంగీతం} mean? (\textquotedblleft{}Music\textquotedblright{}.) Say it once more.
\end{tcolorbox}
\section[\texorpdfstring{\te{నాటకం} (nāṭakaṁ)}{నాటకం (nāṭakaṁ)}]{\te{నాటకం} (nāṭakaṁ) — a play, a drama}

\label{lesson:TE-C275-natakam}



\begin{quote}
Before the new one: say the Telugu for a poet, then the Telugu for music.
\end{quote}

\subsection*{You'll want to know: \te{నాటకం}}
\textbf{\te{నాటకం}} — \emph{nāṭakaṁ} — \textquotedblleft{}a play, a drama\textquotedblright{}.

\begin{grammarlens}[title={What you've built}]
One more word for this chapter.
\end{grammarlens}

\subsection*{Your turn}
\begin{itemize}
\raggedright
  \item \emph{Say it:} \emph{nāṭakaṁ}
  \item \emph{Say it:} \emph{nāṭakaṁ}, once more
  \item \emph{Say it:} say \emph{nāṭakaṁ}
  \item \emph{Recall:} say the Telugu for a poet, then the Telugu for music, then say \emph{nāṭakaṁ} again
\end{itemize}

\begin{tcolorbox}[breakable,colback=teal!4,colframe=teal!35!black,title={Before you move on}]
What does \te{నాటకం} mean? (\textquotedblleft{}A play, a drama\textquotedblright{}.) Say it once more.
\end{tcolorbox}
